\ifdefined\LFSWebBook\else
\documentclass[../textbook.tex]{subfiles}
\begin{document}
\fi
\bookchapter{预训练原理}{ch:pretraining-scaling}

预训练决定模型首先学习怎样的文本分布，以及有限资源应投入多少参数和多少训练词元。两者需要一起考虑：同样的词元数量可能对应不同语言、领域与重复程度；同样的参数数量也可能因训练不足而不能充分发挥作用。规模定律提供的是受条件约束的预算模型，训练目标则规定该模型究竟在衡量什么。

预训练目标由有效监督词元、文档边界和数据混合共同确定，计算预算则约束模型规模与训练词元数。数据生产、资源预算和参数更新分别见\bookxref{ch:data-engineering}、\bookxref{ch:resource-budget}与\bookxref{ch:training-stability}。

\section{预训练所估计的分布}

\subsection{由联合概率到条件预测}
\term{预训练}{Pretraining}{}通常从随机初始化开始，利用大规模数据学习可迁移的表示与生成分布。已有基座上的\term{继续预训练}{Continued Pretraining}{CPT}沿用相近目标，但初始参数已经包含先前训练分布的信息，把它当作同等数据预算下的从零训练会低估其起点。

设文档 $x=(x_1,\ldots,x_T)$，在给定起始标记与终止规则下，链式法则给出
\begin{equation}
 p_\theta(x)=\prod_{t=1}^{T}p_\theta(x_t\mid x_{<t}),\qquad
 -\log p_\theta(x)=-\sum_{t=1}^{T}\log p_\theta(x_t\mid x_{<t}).
 \label{eq:pt-chain}
\end{equation}
$\theta$ 为可学习参数，$x_{<t}$ 是当前位置之前的词元。该概率分解不要求文档内部各词元独立；相关性正是通过条件分布表达的。训练时使用真实前缀，称为\term{教师强制}{Teacher Forcing}{}。推理时前缀包含模型自身生成的内容，因此训练似然改善与生成质量改善有关，却不是完全相同的判据。

若训练样本来自分布 $P$，在固定上下文定义下，期望负对数似然可写为
\begin{equation}
 \mathbb E_{c\sim P}\mathbb E_{y\sim P(\cdot\mid c)}[-\log p_\theta(y\mid c)]
 =\mathbb E_c H(P(\cdot\mid c))+
 \mathbb E_c D_{\mathrm{KL}}(P(\cdot\mid c)\Vert p_\theta(\cdot\mid c)).
 \label{eq:pt-risk}
\end{equation}
这里 $c$ 表示可见上下文，$y$ 为预测目标，$H$ 为条件熵，$D_{\mathrm{KL}}$ 为相对熵。第一项在分布固定时与参数无关，最小化损失因而推动模型逼近该训练条件分布。改变文档筛选、采样权重或可见前缀，就改变了这里的 $P$；数据处理因而直接影响学习目标。

\subsection{符号与假设}
\begin{table}[htbp]
\centering
\caption{预训练目标与规模法则主要符号与计量对象}
\label{tab:rev-pretraining-symbols}
\begin{tabularx}{\linewidth}{lX}
\toprule 符号 & 含义\\\midrule
$P_i,\pi_i$ & 来源 $i$ 的有效目标词元分布及其训练混合权重。\\
$B,T,V$ & 批次序列数、物理块长度与词表大小。\\
$m_{b,t},S$ & 预测位置的监督掩码；当前聚合范围内有效目标总数。\\
$N,D,C$ & 模型参数规模、训练消费词元数与训练浮点运算量。\\
$D_u,D_s,D_p$ & 不同内容的可用词元量、有效监督量、实际处理位置量。\\
$E,A,B_s,\alpha,\beta$ & 经验损失底座、两个正系数及两个正幂律指数。\\
$\kappa,C_0$ & 计算近似中的比例系数与给定计算预算。\\
\bottomrule
\end{tabularx}
\end{table}
$B_s$ 使用下标，以避免与批次大小 $B$ 混淆；$N$ 在本章专指参数规模，不用于表示有效词元数。

\begin{assumption}[预算模型的适用条件]
似然比较固定分词器、验证分布、边界策略与监督区域。规模拟合针对相近的密集 Transformer 架构、数据质量、训练长度与优化完成程度；参数计数和词元计数采用同一口径。幂律描述相应观测范围内的经验关系，其适用范围由拟合数据和训练条件决定。
\end{assumption}

\section{有效词元上的目标与归一化}

\subsection{标签配对与稳定交叉熵}
考虑长度为 $T$ 的输入块 $(x_0,\ldots,x_{T-1})$。模型输出 $\mat{Z}\in\mathbb R^{B\times T\times V}$，位置 $t$ 的向量 $\vect{z}_{b,t,:}$ 预测 $x_{b,t+1}$，因此块内候选目标只有 $T-1$ 个。若库中标签最初复制输入，损失内部须将输出的前 $T-1$ 位与标签的后 $T-1$ 位配对。另一种等价接口显式构造目标序列，并让输入相对于目标右移。两种实现同时移位会得到错误的配对。

令 $m_{b,t}=1$ 表示这一预测配对参与监督，$m_{b,t}=0$ 表示忽略，则
\begin{align}
 \ell_{b,t}&=\log\sum_{v=1}^{V}\exp z_{b,t,v}-z_{b,t,x_{b,t+1}},\\
 S&=\sum_{b=1}^{B}\sum_{t=0}^{T-2}m_{b,t},\qquad
 \mathcal L=\frac{1}{S}\sum_{b,t}m_{b,t}\ell_{b,t}.
 \label{eq:pt-masked-loss}
\end{align}
实现对数求和指数时可先减去最大分数再加回，以避免指数溢出。$S=0$ 的批次没有定义平均训练目标，必须跳过更新或在构造阶段排除，把分母替换成一之后声称取得零损失只是在掩盖这一点。

由 softmax 的导数可得
\begin{equation}
 \frac{\partial\mathcal L}{\partial z_{b,t,v}}
 =\frac{m_{b,t}}{S}\left(p_{b,t,v}-\mathbf1[v=x_{b,t+1}]\right).
 \label{eq:pt-logit-gradient}
\end{equation}
因此，分母不仅改变日志数值，也改变梯度权重。填充掩码控制哪些位置是实际输入，因果掩码控制可见性，监督掩码控制哪些配对贡献梯度；三个对象需要分别定义。

\subsection{批次均值的聚合偏差}
设两个微批次分别含 $S_1=2$、$S_2=6$ 个有效目标，对应平均损失 $L_1=1$、$L_2=3$。整体目标为
\begin{equation}
 L=\frac{2\times1+6\times3}{2+6}=2.5,
\end{equation}
而简单平均得到2。后者给短批次中的每个词元三倍权重。对梯度累积，正确的整体梯度为
\begin{equation}
 g=\sum_k\frac{S_k}{\sum_jS_j}\nabla_\theta L_k.
 \label{eq:pt-accum}
\end{equation}
混合精度和裁剪见\bookxref{ch:training-stability}，分布式梯度平均见\bookxref{ch:distributed-training}；目标层面的原则是先确定总监督量，再分配各微批次的权重。

按文档等权也是合法但不同的目标。若文档 $i$ 的有效长度为 $T_i$，词元等权与文档等权分别为
\begin{equation}
 L_{\mathrm{token}}=\frac{\sum_iT_iL_i}{\sum_iT_i},\qquad
 L_{\mathrm{doc}}=\frac1M\sum_{i=1}^{M}L_i.
\end{equation}
文档等权提高短文档中单个词元的影响。选择哪一种取决于希望拟合的统计单位，由数据加载器的默认归约决定则属于偶然。

\subsection{困惑度的比较条件}
\term{困惑度}{Perplexity}{PPL}定义为 $\exp(\mathcal L)$，其中对数以自然底计算。它等于正确目标概率几何平均数的倒数。上例整体困惑度为 $\exp(2.5)\approx12.18$，它既不是两个批次困惑度的算术平均，也不是 $\exp(2)$。

困惑度的含义是模型可能输出的词数或事实错误率，这两种常见理解都与定义不符。它的比较依赖相同分词器、验证文本、可见上下文和目标掩码。使用相同文本而改变分词器，会改变目标单位及条件分解；此时可以补充每字节负对数似然，但仍须统一规范化和文本覆盖。\footnote{以自然对数计量的损失常称为每词元 nat；换算为 bit 应除以 $\log 2$。改变对数底不改变同一实验中的排序，却会改变所报告的系数与数值。}

\section{序列装箱的条件分布}

\subsection{终止标记与注意力隔离}
\term{序列装箱}{Sequence Packing}{}把多个片段安排进物理块以减少空位。若独立文档 $a$、$b$ 被连成 $a,\mathrm{EOS},b,\mathrm{EOS}$，采用普通因果注意力时，$b$ 的预测条件包含 $a$。终止标记提供边界信息，独立文档目标还需阻止跨文档读取。

若训练目标要求文档独立，设位置 $t$ 的文档编号为 $d(t)$，可见性需要满足
\begin{equation}
 A_{t,j}=\mathbf1[j\le t]\mathbf1[d(j)=d(t)].
 \label{eq:pt-document-mask}
\end{equation}
同时必须处理跨段目标配对。屏蔽 $a$ 对 $b$ 的注意力，并不会阻止 $a$ 的末尾输出预测 $b$ 的第一个词元。通常每段含起始标记，由起始标记预测该段首词元；上一段终止位置到下一段起始位置的配对不计入损失。是否重置位置编号也需固定，因为位置分布会因此改变。

\begin{figure}[htbp]
\centering
\begin{tikzpicture}[x=1.45cm,y=.9cm,>=stealth,every node/.style={font=\small}]
\foreach \x/\txt in {0/BOS,1/$a_1$,2/EOS,3/BOS,4/$b_1$,5/EOS}{
 \node[draw,rounded corners,minimum width=1.15cm,minimum height=.55cm] (p\x) at (\x,0) {\txt};}
\draw[->] (p0.south) to[bend right=35] node[below]{监督} (p1.south);
\draw[->] (p1.north) to[bend left=35] (p2.north);
\draw[dashed,->] (p2.south) to[bend right=35] node[below]{忽略} (p3.south);
\draw[->] (p3.north) to[bend left=35] (p4.north);
\draw[->] (p4.south) to[bend right=35] node[below]{监督} (p5.south);
\draw[densely dotted] (2.5,-1.05)--(2.5,.85);
\node at (1,-1.3){文档 $a$ 内因果可见};\node at (4,-1.3){文档 $b$ 内因果可见};
\end{tikzpicture}
\caption{独立文档装箱的两种边界：注意力禁止跨文档读取，损失忽略跨文档预测配对。实线表示有效相邻监督。}
\label{fig:pt-packing}
\end{figure}

对于连续章节、代码文件集合或保留上下文的对话，跨片段可见也可能是刻意设计。关键在于物理切片是否对应统计独立单位。把同一文档拆段后完全隔离，会损失长距离条件；把不相关文档连接而不隔离，则引入人为前缀分布。

\subsection{块边界与尾部损失}
连续词元流切成多个独立输入块时，每个块最后位置缺少块内后继目标。若采用重叠一个词元的窗口，可保留该边界预测，但更大的重叠会重复监督内容。若只保留重叠区域作为上下文，应将已经计入过的目标屏蔽，以免改变权重。

尾部丢弃、填充、跨分片续接分别影响覆盖率、计算利用率和读取状态。假设100个文件各有129个词元，每个文件单独按128切块会丢弃100个词元；跨文件保留缓冲尾部可减少这一损失，却必须另外保留文档边界，以减少浪费为由改变条件结构会使目标偏离原设定。验证集也必须使用同一明确策略，否则损失变化可能只是可见上下文或被评分目标不同。

\begin{note}[装箱率与监督率是不同指标]
非填充位置占物理位置的比例描述装箱利用率；有效目标占物理位置的比例描述监督密度。边界标记、仅供上下文的前缀和掩码区域均可能消耗计算而不贡献直接损失。规模预算必须同时保留这两类计数。
\end{note}

\section{数据混合与长尾覆盖}

\subsection{重要性权重}
设来源 $i$ 的目标词元分布为 $P_i$，希望以权重 $\pi_i$ 训练，则
\begin{equation}
 P_{\mathrm{train}}=\sum_i\pi_iP_i,\qquad
 L(\theta)=\sum_i\pi_iL_i(\theta),\qquad \sum_i\pi_i=1.
 \label{eq:pt-mixture}
\end{equation}
若加载器先以概率 $q_i$ 选择来源，再抽取平均有 $\bar s_i$ 个有效目标的文档，聚合后的词元比例近似为
\begin{equation}
 \pi_i=\frac{q_i\bar s_i}{\sum_jq_j\bar s_j}.
 \label{eq:pt-source-length}
\end{equation}
因此，两个来源按文档各取一半，若平均有效长度分别为100与900，实际词元比例为一比九。要得到词元各半，需要文档采样概率九比一。对具有不同监督密度的任务，应使用有效目标长度而不只使用原始文档长度。

常用温度平滑以合格可用词元量 $n_i$ 定义 $\pi_i\propto n_i^\tau$，$0\le\tau\le1$。这里 $\tau$ 是混合指数，并非 softmax 温度或规模幂律指数。$\tau=1$ 保留规模比例，$\tau=0$ 对非空来源等权。降低指数会提高小来源相对份额，该来源实际拥有的不同信息则不因此增加。

\subsection{数据混合配比}
\begin{example}


设三个来源可用词元量为90、9、1十亿。自然比例为 $0.90,0.09,0.01$；取 $\tau=\frac{1}{2}$ 后，权重与 $\sqrt{90},3,1$ 成正比，约为 $0.7034,0.2224,0.0741$，最后一项因四舍五入可作相应校正。计划消费200十亿词元时，各来源的名义遍历倍数为
\begin{equation}
 e_i=\frac{D\pi_i}{n_i},
\end{equation}
约为 $1.56,4.94,14.83$。小来源的词元权重提高约七倍，其重复程度也显著增加。该指标衡量总体消费量与库存量之比，各文档访问次数由采样策略决定。

带替换采样 $m$ 次、均匀抽取 $n$ 个不同对象时，期望见过的不同对象数为
\begin{equation}
 U=n\left[1-\left(1-\frac1n\right)^m\right].
 \label{eq:pt-unique}
\end{equation}
推导只需对每个对象定义是否至少出现一次的指示变量，再求期望之和。$m=n$ 时，大样本近似只有 $(1-\conste{}^{-1})n$ 个不同对象被访问。真实文档长度不等且存在近重复，此对象公式只能给出不同信息词元数的一个上界。
\end{example}

\subsection{训练分布与评测分布}
训练混合比例表达资源投入，评测混合比例表达质量偏好，两者不必相同。若目标评测分布权重为 $r_i$，训练采样权重为 $\pi_i>0$，可用 $\frac{r_i}{\pi_i}$ 对来源损失作重要性加权，使期望匹配目标，但极端权重会放大方差。若目标来源完全没有训练支持，加权也产生不出缺失的样本。

质量筛选也会改变语言和文体覆盖。低困惑度文本可能更贴近评分模型熟悉的分布，却未必覆盖罕见术语或形式化推导。预训练配方应分别记录准入、质量层、语言、领域、长度和重复程度；压成一个“优质分数”会丢掉这些可区分的维度。数据流水线的实现见\bookxref{ch:data-engineering}，本章关心这些决策怎样改变 $P_{\mathrm{train}}$。

\subsection{课程与继续预训练}
\term{课程学习}{Curriculum Learning}{}使混合权重或难度随训练进程变化，记为 $\pi_i(u)$，其中 $u$ 可以是累计词元数。即便总消费量相同，先通用后领域与交错混合也可能得到不同参数，因为连续梯度更新不可交换。课程应同时说明阶段长度、过渡方式和通用能力保持条件。

继续预训练通常需要评估领域收益与旧分布退化。保留部分通用数据可以约束漂移，却没有一个跨模型通用的保证比例。小规模配方比较应固定有效词元预算、候选模型与评测切片，再决定是否扩展；测试集不得反复用于选择混合权重。

\section{不同预训练任务的监督结构}

自回归、掩码恢复与去噪的区别首先在于可见信息和被预测变量，模型名称只是次要线索。统一表示为：从原始文档 $x$ 通过变换 $q(c,y\mid x)$ 产生条件 $c$ 与目标 $y$，优化
\begin{equation}
 L_q(\theta)=\mathbb E_{x}\mathbb E_{(c,y)\sim q(\cdot\mid x)}
 \left[-\sum_t\log p_\theta(y_t\mid c,y_{<t})\right],
 \label{eq:pt-corruption}
\end{equation}
并明确按目标词元还是样本归一化。对于并行预测多个掩码位置的模型，求和项使用各位置的条件概率，其目标描述被遮蔽位置的条件预测。

\begin{table}[htbp]
\centering
\caption{不同自监督预训练任务的条件结构与监督边界对比}
\label{tab:rev-pretraining-tasks}
\begin{tabularx}{\linewidth}{lXXX}
\toprule 任务 & 可见信息 & 预测内容 & 主要边界\\\midrule
自回归语言建模 & 当前词元之前的前缀 & 下一词元 & 直接匹配逐步续写；训练前缀是真实数据。\\
掩码语言建模 & 含遮蔽的双向上下文 & 被选中的词元 & 擅长双向表示；训练损失与原文自回归困惑度属于不同口径。\\
跨度去噪 & 删除跨度后的文档与哨兵标记 & 被删跨度的序列 & 需要固定破坏分布、跨度顺序和停止协议。\\
前缀语言建模 & 完整前缀与已生成后缀 & 后续部分 & 前缀通常不计目标损失，监督密度随切分点变化。\\
中间填充 & 原文前段与后段 & 中间缺失内容 & 需要重排及边界标记；序列顺序不同于原文。\\
\bottomrule
\end{tabularx}
\end{table}

T5 的研究比较了文本到文本框架中的去噪目标与迁移设置\cite{raffel2020t5}。以原文“甲 乙 丙 丁 戊”为例，删除“乙 丙”并以哨兵 $s_0$ 替代，条件可写为“甲 $s_0$ 丁 戊”，目标为“$s_0$ 乙 丙 $s_1$”；哨兵用于定位缺失跨度，按通用起始标记解释它会得到错误的条件结构。这里的符号例用于说明条件结构，具体制品的终止标记仍以其协议为准。

\term{中间填充}{Fill-in-the-Middle}{FIM}可以将文档拆为前段 $a$、中段 $b$、后段 $c$，重排为带边界标记的 $a,c,b$，从而让因果模型在生成 $b$ 时同时读取原文两侧\cite{bavarian2022fim}。它保持因果注意力只见重排序列的前缀，通过改变顺序使原文两侧内容都成为已知前缀。

混合任务可写为 $L=\sum_k\lambda_kL_k$，但 $\lambda_k$ 的含义依赖 $L_k$ 的归一化。如果一个任务只监督15\%位置，另一个监督近乎全部位置，相同样本抽样率对应的目标词元贡献相差数倍。比较任务还应报告处理位置量和实际计算量，否则较低损失可能来自更容易的条件或更高的监督资源，其中并不必然包含更优的表示学习机制。

\section{参数、数据与计算的统一预算}

\subsection{词元计量口径}
$D_u$ 表示去重后可用内容的规模，$D_s$ 表示累计有效目标数，$D_p$ 表示模型实际处理的位置数。重复多个遍历增加 $D_s$ 和 $D_p$，不会同比增加 $D_u$；填充和仅供上下文的位置增加 $D_p$，不一定增加 $D_s$。

设每次参数更新有 $R$ 个数据并行进程，每进程累积 $G$ 个微批次，每批次 $B$ 个长度 $T$ 的块，训练 $K$ 步，在规则相同且无额外掩码时
\begin{equation}
 D_p=KRGBT,\qquad D_s=KRGB(T-1).
 \label{eq:pt-budget-count}
\end{equation}
例如 $K=1000,R=4,G=2,B=8,T=1024$，得到 $D_p=65536000$，$D_s=65472000$。跨段屏蔽和填充会进一步降低 $D_s$，实际应累加监督掩码，$T-1$ 只在边界规则相同时成立。

\subsection{训练计算近似的来源}
对主要由密集矩阵乘法构成的网络，前向每个词元约需与参数规模成比例的运算；反向包含输入与权重梯度两部分。因此常用
\begin{equation}
 C\approx\kappa ND,\qquad \kappa\approx6
 \label{eq:pt-compute}
\end{equation}
作为数量级近似。这里乘加计为两次浮点运算，$D$ 通常接近处理词元量；只有在监督密度高且边界损失很小时，才近似等于有效目标量。\footnote{参数口径可能包含或排除嵌入，稀疏专家模型还须区分总参数与激活参数。引用某篇规模拟合的系数时，更换参数定义就不能沿用原数值。}

式\eqref{eq:pt-compute}没有显式表达随上下文增长的注意力开销、激活重计算和特殊算子成本。长上下文时应重新估计 $\kappa$ 或改用更详细模型。设备峰值乘运行时间是硬件上限；通信、存储等待与内核效率决定实际完成的模型计算，峰值预算与损失定律之间还隔着这段差距。

\section{经验规模定律的拟合}

\subsection{幂律描述与数据条件}
\term{规模定律}{Scaling Law}{}以经验函数描述模型规模、数据和计算与损失之间的关系。Kaplan 等人的研究系统分析了语言模型交叉熵随这些变量变化的规律\cite{kaplan2020scaling}；后续 Chinchilla 工作重新研究固定训练计算下的参数与训练词元分配\cite{hoffmann2022chinchilla}。这些研究的配置与拟合程序不同，把其中的指数拼接成一条统一物理定律会掩盖这一差异。

本章使用一个便于分析的加性模型
\begin{equation}
 L(N,D)=E+A N^{-\alpha}+B_sD^{-\beta},\qquad
 A,B_s,\alpha,\beta>0.
 \label{eq:pt-scaling}
\end{equation}
$E$ 表示该经验模型中的渐近底座，$AN^{-\alpha}$ 与 $B_sD^{-\beta}$ 分别描述参数和数据受限时的损失余量。数据、分词器和条件窗口的变化也会改变渐近损失。表\ref{tab:rev-scaling-laws-comparison}对比了代表性规模法则拟合方案的经验模型形式、拟合指数与工程结论。

\begin{table}[htbp]
\centering
\small
\setlength{\tabcolsep}{3.5pt}
\caption{主流语言模型预训练规模法则（Scaling Laws）参数与拟合结论对比}
\label{tab:rev-scaling-laws-comparison}
\begin{tabularx}{\linewidth}{lp{28mm}cp{24mm}X}
\toprule
拟合方案 & 损失经验模型 & 指数 $(\alpha, \beta)$ & 最优分配关系 & 核心差异与工程意义 \\
\midrule
Kaplan 等 (2020) & $L(N,D)=E+(\frac{N_c}{N})^\alpha+(\frac{D_c}{D})^\beta$ & $(0.076, 0.057)$ & $N \propto C^{0.73},\, D \propto C^{0.27}$ & 倾向模型增大快于数据规模，低估了数据量对计算最优的贡献 \\
Chinchilla (2022) & $L(N,D)=E+\frac{A}{N^\alpha}+\frac{B_s}{D^\beta}$ & $(0.34, 0.28)$ & $N \propto C^{0.5},\, D \propto C^{0.5}$ & 修正余弦退火协议；证明参数量与词元数应等比例协同扩展 \\
Llama 3 (2024) & 包含退火与二阶对数外推 & -- & 推理主导倾斜 ($D \ge 100N$) & 考虑推理阶段服务成本的长期分摊，在固定模型规模下进行深度超训 \\
\bottomrule
\end{tabularx}
\end{table}

单独固定数据研究参数增长时，若已知底座 $L_\infty$，有
\begin{equation}
 \log(L-L_\infty)=\log A-\alpha\log N.
\end{equation}
例如参数每扩大四倍、损失余量减半，则 $\alpha=\frac{\log2}{\log4}=\frac{1}{2}$。若对 $\log L$ 而非 $\log(L-L_\infty)$ 作线性拟合，底座占比较高时斜率会明显偏小。把未知底座与指数同时从狭窄区间估计，也容易产生多组近似等价解。

\subsection{拟合应覆盖两个维度}
仅在一条 $ND$ 固定曲线上取样，会使参数与数据强相关，难以分别识别两者贡献。应在多组参数规模和训练长度上设计实验，并保留独立规模点检验预测。每个实验须具有可比的优化完成程度：若大模型使用明显不合适的学习率，拟合可能把优化不足误归因为规模效率差。

一种拟合形式为
\begin{equation}
 \min_{E,A,B_s,\alpha,\beta}\sum_{j}w_j\,
\rho\!\left(
 \log L_j-\log[E+AN_j^{-\alpha}+B_sD_j^{-\beta}]\right),
 \label{eq:pt-fit}
\end{equation}
其中 $j$ 为实验点，$w_j$ 为可靠性权重，$\rho$ 可取平方损失或稳健损失。使用对数残差表示偏重相对误差；误差形式应与预算决策的目标相匹配。约束正系数可通过指数参数化实现。

同一训练运行中的多个检查点共享数据与随机轨迹，把它们当成相互独立实验会人为放大置信度。应按训练运行或随机种子组织重采样，并报告外推点的误差区间。残差若随规模、长度或来源系统变化，说明模型遗漏了因素，增加小数位并不能提高可信度。

\begin{algorithm}[规模拟合与预算候选生成]
\begin{enumerate}
\item 冻结分词器、数据混合、评测集合、上下文和参数计数口径；保留测试集不参与配方选择。
\item 设计覆盖多个 $N$ 与 $D$ 的训练网格，为每个规模分配合理优化配置，记录真实消费量与计算估计。
\item 以有效目标总和聚合固定验证集损失，保留运行身份、重复程度与异常说明。
\item 拟合式\eqref{eq:pt-scaling}，在未参与拟合的规模点比较预测与实际损失，分析残差及参数不确定性。
\item 在给定预算下求连续候选，再映射到可部署的离散架构与可用数据范围；对相邻候选开展预算内比较。
\item 保存数据配方、拟合区间、候选选择依据与中止条件，使后续扩大规模有明确依据。
\end{enumerate}
\end{algorithm}

\section{计算最优分配的推导}

\subsection{拉格朗日条件}
将式\eqref{eq:pt-scaling}与 $C=\kappa ND\le C_0$ 联立。因为在假设范围内损失随 $N,D$ 单调下降，内部最优点会用尽预算。定义拉格朗日函数
\begin{equation}
 \mathcal J=E+AN^{-\alpha}+B_sD^{-\beta}
 +\lambda(\kappa ND-C_0).
\end{equation}
分别对 $N,D$ 求导并令其为零：
\begin{align}
 -\alpha AN^{-\alpha-1}+\lambda\kappa D&=0,\\
 -\beta B_sD^{-\beta-1}+\lambda\kappa N&=0.
\end{align}
第一式乘 $N$，第二式乘 $D$，消去相同的 $\lambda\kappa ND$，得到
\begin{equation}
 \alpha AN^{-\alpha}=\beta B_sD^{-\beta}.
 \label{eq:pt-balance}
\end{equation}
该条件平衡的是按对数规模计算的边际收益；只有 $\alpha=\beta$ 时，两项损失余量才相等。

令 $K_c=\frac{C_0}{\kappa}$，代入 $D=\frac{K_c}{N}$，可得
\begin{align}
 N_*&=\left(\frac{\alpha A}{\beta B_s}\right)^{\frac{1}{\alpha+\beta}}
 K_c^{\frac{\beta}{\alpha+\beta}},\\
 D_*&=\left(\frac{\beta B_s}{\alpha A}\right)^{\frac{1}{\alpha+\beta}}
 K_c^{\frac{\alpha}{\alpha+\beta}}.
 \label{eq:pt-optimal}
\end{align}
最优解由所选经验函数和计算约束共同确定，应用前需核对拟合范围与预测误差。

\subsection{最优点的唯一性与预算指数}
令 $u=\log N$，固定预算后的非底座损失为
\begin{equation}
 f(u)=A\conste{}^{-\alpha u}+B_sK_c^{-\beta}\conste{}^{\beta u}.
\end{equation}
其二阶导数
\begin{equation}
 f''(u)=\alpha^2A\conste{}^{-\alpha u}+\beta^2B_sK_c^{-\beta}\conste{}^{\beta u}>0.
\end{equation}
因此在没有额外边界的连续正数域中，驻点是唯一全局最优点。存在显存上限、最小模型规模或数据库存约束时，最优解可能落在可行区边界。

最优规模随预算的指数分别为 $\frac{\beta}{\alpha+\beta}$ 与 $\frac{\alpha}{\alpha+\beta}$。当两指数相等时，模型和数据都按预算平方根增加；预算扩大四倍时，两者各扩大两倍。最优余量则满足
\begin{equation}
 L_*(C)-E\propto C^{-\frac{\alpha\beta}{\alpha+\beta}}.
\end{equation}
收益随计算增长而递减。若 $\alpha\ne\beta$，最优 $\frac{D_*}{N_*}$ 也随预算变化；最佳词元参数比应按式\eqref{eq:pt-optimal}计算。

\subsection{计算最优分配}
\begin{example}


例如，定义 $n=\frac{N}{10^9}$、$d=\frac{D}{10^9}$，假设
\begin{equation}
 L(n,d)=2+n^{-\frac{1}{2}}+4d^{-\frac{1}{2}},\qquad C_0=\num{9.6e19},\quad\kappa=6.
\end{equation}
计算预算给出 $nd=16$。由边际平衡 $n^{-\frac{1}{2}}=4d^{-\frac{1}{2}}$，得 $d=16n$，联立后 $n_*=1,d_*=16$，预测损失为4。

\begin{table}[htbp]
\centering
\caption{固定计算预算 $C_0$ 下不同参数规模与词元数量分配的预测损失对比}
\label{tab:rev-isocompute-allocation}
\begin{tabular}{rrrr}
\toprule 参数 $n$ & 词元 $d$ & 乘积 $nd$ & 预测损失\\\midrule
0.25 & 64 & 16 & 4.5\\
1 & 16 & 16 & 4.0\\
4 & 4 & 16 & 4.5\\
\bottomrule
\end{tabular}
\end{table}

较小模型获得更多数据但容量项较大，较大模型压缩了训练长度而数据项较大。两端计算量相同，参数规模本身并不能决定结果。预算扩大四倍后，最优为 $n=2,d=32$，损失 $2+\sqrt2\approx3.414$。

\begin{figure}[htbp]
\centering
\begin{tikzpicture}[x=2.4cm,y=1.2cm,>=stealth,every node/.style={font=\small}]
\draw[->] (-1.45,0)--(1.65,0) node[right]{$\log_4 n$};
\draw[->] (-1.35,-.05)--(-1.35,3.1) node[above]{$L-2$};
\draw[thick] plot[domain=-1:1,samples=60] (\x,{pow(4,-\x/2)+pow(4,\x/2)});
\foreach \x/\y/\lbl in {-1/2.5/{0.25},0/2/{1},1/2.5/{4}}{
 \fill (\x,\y) circle (2pt);\node[below] at (\x,0){$\lbl$};}
\node[below] at (0,-.4){刻度文字为参数规模 $n$};
\node[above] at (0,2){计算固定，$d=\frac{16}{n}$};
\node[left] at (-1.35,2){2};\node[left] at (-1.35,2.5){2.5};
\end{tikzpicture}
\caption{数值例在固定计算下的损失余量。曲线由本节解析函数计算；纵轴从零起算，曲线最低点对应参数与数据的边际收益平衡。}
\label{fig:pt-isocompute}
\end{figure}

若可用不同数据只有四十亿词元且不允许重复，则 $d\le4$。在相同预算下，连续可行最优落在 $d=4,n=4$，损失为4.5。重复训练改变有效数据量与损失的关系，需要根据重复次数重新拟合。
\end{example}

\section{规模结论的外推边界}

\subsection{数据受限与优化受限}
经验式中的数据量通常建立在特定重复范围和质量分布上。重复曝光有时仍能改善优化，可提供的新统计信息却是有限的。合成数据也必须考察新信息、错误相关性与验证覆盖，文件数增加不等于原始独立数据预算同步增加。

同样，计算量相同但学习率、批次、衰减与停止时机不同，可能导致不同损失。规模最优讨论要求各候选达到可比的预算内训练效率，所有模型机械使用同一更新步数并不满足这一要求。\bookxref{ch:training-stability}将展开优化器状态与调度；在本章的预算实验中，它们属于必须记录的控制条件。

\subsection{训练最优与全生命周期最优}
Chinchilla 式问题关注给定训练计算的模型选择。部署时还可能关注模型容量、延迟和累计推理量。若未来生成 $Q$ 个词元，简化生命周期成本可写为
\begin{equation}
 C_{\mathrm{life}}\approx\kappa ND+\kappa_{\mathrm{inf}}NQ,
\end{equation}
其中 $\kappa_{\mathrm{inf}}$ 是给定推理场景下的近似系数。较小模型经过更多训练，可能在大量调用时降低总成本；这也改变了约束与优化目标，仅训练计算最优点在此不再适用。真实推理成本还随上下文、缓存与硬件变化，上式说明的是新增目标项的方向。

\subsection{损失与能力的不同尺度}
验证交叉熵衡量预训练分布上的平均预测质量；下游任务还受评分阈值、提示方式、样本量和数据混合影响。模型选择应同时保留可比的预训练损失与领域、语言、长上下文等任务指标，分别观察总体拟合与能力切片。

\begin{note}[规模定律的正确用法]
规模定律用于缩小预算候选范围，并提出可被实验否定的预测。拟合区间、分布、架构或目标发生变化后，应重新评估系数与误差。连续最优解是架构选型的起点，最终方案还须满足数据可用量、设备容量和部署要求。
\end{note}

\section{预训练的数据流与状态流}

\begin{figure}[htbp]
\centering
\begin{tikzpicture}[>=stealth,node distance=7mm,every node/.style={font=\small},box/.style={draw,rounded corners,align=center,text width=28mm,minimum height=10mm}]
\node[box] (data) {版本化来源\\混合与边界};
\node[box,right=of data] (batch) {分词与装箱\\$x,m,A$};
\node[box,right=of batch] (model) {模型前向\\$Z[B,T,V]$};
\node[box,right=of model] (loss) {目标配对\\有效词元损失};
\node[box,below=12mm of loss] (update) {反向与更新\\实际训练计数};
\node[box,left=of update] (state) {完整检查点\\参数与消费状态};
\node[box,left=of state] (eval) {固定验证分布\\规模与配方选择};
\draw[->] (data)--(batch);\draw[->] (batch)--(model);\draw[->] (model)--(loss);
\draw[->] (loss)--(update);\draw[->] (update)--(state);\draw[->] (state)--(eval);
\draw[->] (update.east)--++(.3,0)|-(model.north);
\draw[dashed,->] (state)--(batch);\draw[dashed,->] (eval)--(data);
\end{tikzpicture}
\caption{预训练同时积累模型知识与可恢复状态。实线表示计算和保存路径，虚线表示恢复或下一轮实验设计；验证结果不直接生成当前训练目标。}
\label{fig:pt-system}
\end{figure}

一次训练状态可记为 $s_k=(\theta_k,o_k,h_k,r_k,d_k)$，分别表示参数、优化器状态、调度状态、随机状态与数据消费位置。一次更新为 $s_{k+1}=F(s_k,\mathcal B_k)$。只恢复参数会改变后续状态转移，这样得到的续训与原始轨迹并不等价。数据配比随进度变化时，还需要恢复其阶段位置；流式洗牌则需要能够重建缓冲与分片消费关系。

\begin{algorithm}[按有效词元推进预训练]
\begin{enumerate}
\item 固定初始模型、分词器、来源版本、配比日程、边界策略和总预算，分别初始化处理量与有效监督量计数。
\item 从可恢复采样状态读取下一组文档，构造输入、段编号、注意力与监督掩码；保留来源跨度以归属消费量。
\item 对一个更新范围汇总有效目标数，按式\eqref{eq:pt-masked-loss}与式\eqref{eq:pt-accum}形成正确归一化梯度；无有效目标时不推进参数更新。
\item 完成优化器更新后推进调度和数据计数；记录各来源实际有效词元份额、重复程度与计算消耗。
\item 在预定评估点累计固定验证集的负对数似然总和与有效目标总数，分别报告总体及来源切片。
\item 在明确更新边界保存完整状态及不可变数据引用；达到预算后交付所选检查点、计量口径与独立评测结果。
\end{enumerate}
\end{algorithm}

训练过程平均损失反映历史参数在各批次上的表现，与最终参数的验证损失是两个量。用于选择最佳检查点的验证集已经参与模型选择，最终测试集应继续隔离。




\chapterreferences
\lfsendchapterdocument
